\subsection{Word Segmentation}
\label{sec:word-segmentation}
\slidesrc{Syntax: RE \& word: 16--20}

\begin{nbdef}{Word segmentation}
\textbf{Word segmentation} (also \emph{tokenisation}) is the problem of dividing a string
of written language into its component words.
\end{nbdef}

\subsubsection{Why it is not trivial}
\slidesrc{Syntax: RE \& word: 17}

In English and in many other languages written with the Latin alphabet, the
\textbf{space} is a good approximation of a word divider (word delimiter). It is only an
approximation, because languages vary in how they treat collocations and compounds
(Section~\ref{sec:morphology}): many English compound nouns are written in several ways,
e.g.\ \emph{multi-agent} = \emph{multiagent} = \emph{multi agent}, and speakers differ on
whether they think of them as one noun or as a noun phrase. Punctuation adds more cases
(\emph{don't}, \emph{U.S.A.}, \emph{3.14}, \emph{e-mail}).

Many scripts have no word space at all, and there word segmentation is a hard problem:
\begin{itemize}[nosep]
  \item \textbf{Chinese} and \textbf{Japanese}: sentences are delimited, words are not;
  \item \textbf{Thai} and \textbf{Lao}: phrases and sentences are delimited, words are not;
  \item \textbf{Vietnamese}: syllables are delimited, words are not.
\end{itemize}

\subsubsection{The MaxMatch algorithm}
\slidesrc{Syntax: RE \& word: 18}

A simple baseline for Chinese is the greedy \textbf{maximum matching} algorithm. It needs
a \textbf{dictionary} (word list) of the language.

\begin{nbdef}{MaxMatch}
\begin{enumerate}[nosep]
  \item Put a pointer at the beginning of the string.
  \item Choose the \textbf{longest} dictionary word that matches the input starting at
        the pointer, and move the pointer to the end of that word.
  \item If no word matches, move the pointer one character forward instead (creating a
        one-character ``word'').
  \item Repeat from the new pointer position until the end of the string.
\end{enumerate}
\end{nbdef}

\begin{nbexample}{MaxMatch on English without spaces}
Input: \texttt{wecanonlyseeashortdistanceahead} (``we can only see a short distance
ahead'').
\begin{center}
\small
\begin{tabular}{@{}l l l@{}}
\toprule
\textbf{Remaining input} & \textbf{Longest word} & \textbf{Why} \\
\midrule
\texttt{wecanonly\ldots}  & \emph{we}       & \\
\texttt{canonlysee\ldots} & \emph{canon}    & longer than \emph{can} \\
\texttt{lyseea\ldots}     & \emph{l}        & no word starts here: one character \\
\texttt{yseea\ldots}      & \emph{y}        & same \\
\texttt{seeashort\ldots}  & \emph{see}      & \\
\texttt{ashort\ldots}     & \emph{ash}      & longer than \emph{a} and \emph{as} \\
\texttt{ort\ldots}        & \emph{o}, \emph{r}, \emph{t} & three single characters \\
\texttt{distanceahead}    & \emph{distance}, \emph{ahead} & \\
\bottomrule
\end{tabular}
\end{center}
Output: \emph{we canon l y see ash o r t distance ahead}.
\end{nbexample}

\begin{nbintuition}[why greedy fails here and works for Chinese]
Greedy choices are never undone: once \emph{canon} is taken, the correct \emph{can only}
is impossible. In English, words are long and many short words hide inside longer ones,
so this happens constantly. Chinese words are mostly one or two characters long, so the
longest match is usually right, and MaxMatch is a surprisingly strong baseline there.
Better segmenters score \emph{whole} segmentations (e.g.\ with word probabilities)
instead of committing greedily.
\end{nbintuition}

\subsubsection{NLTK}
\slidesrc{Syntax: RE \& word: 19--20}

The course uses \textbf{NLTK} (Natural Language Toolkit, \url{https://www.nltk.org/})
to show some problems and some of their possible solutions. NLTK is a Python platform
for working with human language data: it provides interfaces to over 50 corpora and
lexical resources such as WordNet (Section~\ref{sec:wordnet}), and libraries for
classification, tokenisation, stemming, tagging, parsing and semantic reasoning. The
course also uses less known, but elegant, tools.

\begin{minted}{python}
import nltk
# nltk.download()   # only once, to download corpora and models
from nltk.tokenize import word_tokenize
print(word_tokenize("Don't segment U.S.A. badly, Mr. Smith!"))
# ['Do', "n't", 'segment', 'U.S.A.', 'badly', ',', 'Mr.', 'Smith', '!']
\end{minted}

Note how the tokeniser splits \emph{don't} into \emph{Do} + \emph{n't} (two
morphemes, Section~\ref{sec:morphology}) but keeps \emph{U.S.A.} and \emph{Mr.}
together: the full stop is not always a word boundary.
